\begin{table}[htp]
\centering

    % Esta primera línea define las columnas de la tabla. Los posibles tipos de columna son:
    % c: texto centrado
    % l: texto alineado a la izquierda
    % r: texto centrado a la derecha
    % p: columna de ancho fijo
    % Las columnas tienen ancho dinámico según la anchura máxima de los elementos que contengan.

    % Las columnas l/r/c no parten el texto en filas diferentes si éste es demasiado largo. Para ello, puede utilizar el tipo de columna de ancho fijo "p".
    
    % Las barras verticales | se usan para definir los bordes verticales de la tabla. Pruebe a eliminar algunas y observe qué ocurre.
    \begin{center}
    \begin{tabular}{ | l | l | l | }
        
        % A continuación van las filas de la tabla. En cada fila, las columnas se separan con el carácter &
        % Para terminar una fila se usa \\
        % Para incluir un borde horizontal entre filas se usa \hline
        
        % Cabecera con textos en negrita:
        \hline
        \textbf{Measure} & \textbf{Animation}  & \textbf{Polygon} \\
        \hline
        \textbf{Task clock (miliseconds)}  & 3,561 & 3,574 \\
        \hline
        \textbf{CPUs utilized}  & 0.367 & 0.915 \\
        \hline
        \textbf{Page faults} & 3,609  & 3,554\\
        \hline
        \textbf{Frequency (GHz)} & 3.187 & 3.268 \\
        \hline
        \textbf{Instructions per cycle} & 1.04 & 0.77 \\
        \hline
        \textbf{Branches per second} & 665.660 & 550.479 \\
        \hline
        \textbf{Branch misses percentage} & 1.21 & 1.14 \\
        \hline
        \textbf{Time elapsed (seconds)} & 9.7086 & 3.9050 \\
        \hline
        \textbf{Time user (seconds)} & 2.2427 & 1.1099 \\
        \hline
        \textbf{Time system (seconds)} & 1.3618 & 2.4674 \\
        \hline
        \textbf{Context switches} & 18,666 & 1,914 \\
        \hline
        \textbf{CPU migrations} & 331 & 26 \\
        \hline
        
    \end{tabular} 
    \end{center}
    
    \caption{Comparación de los resultados de la herramienta \textit{perf} durante la ejecución de las pruebas.}
    
    \label{tbl:perf-comparison}
\end{table}